% ============================================================
% ch01_dipole.tex — EXPANDED
% The Cosmic Dipole Anomaly
% Part of: Tetrad-Based Departure Decomposition (Jiwon Park)
% Session S10 draft — 2026-03-30
% ============================================================

\chapter{The Cosmic Dipole Anomaly}
\label{ch:dipole}

The cosmic dipole anomaly is the primary observational
motivation for this work.  Independent datasets spanning
radio continuum, mid-infrared quasars, and peculiar velocity
surveys each report a matter dipole in excess of the FLRW
kinematic prediction, with a combined significance exceeding
$5\sigma$.  We review the observational evidence in
Section~\ref{sec:dipole-obs}, the statistical methodology in
Section~\ref{sec:dipole-stats}, the peculiar velocity context
in Section~\ref{sec:dipole-pv}, the theoretical status of
candidate mechanisms in Section~\ref{sec:dipole-theory}, the
observer-frame discriminator in
Section~\ref{sec:dipole-discriminator}, and the near-term
experimental outlook in
Section~\ref{sec:dipole-outlook}.  The anomaly parameters
relevant to the departure framework are summarised in
Section~\ref{sec:dipole-summary}.

The matter-dipole premise is actively contested.  We therefore treat
CatWISE, radio, and bulk-flow inputs as conditional probes rather than
settled evidence for anisotropic geometry.  Recent reassessments of
CatWISE clustering and mask coupling~\cite{Bashir2025,VonHausegger2025CatWISEClustering},
current radio-source-count analyses~\cite{Boehme2025}, and
tilted-anisotropic model constraints~\cite{Martin2025TiltedAnisotropicDipole}
are discussed as premise and systematics controls.  This framing makes
the likelihood layer updateable: replacing or down-weighting a survey
changes the conditional HTT evidence surface without changing the exact
claim-tier and provenance machinery.


% ============================================================
\section{Number-count dipole excess}
\label{sec:dipole-obs}
% ============================================================

The Cosmological Principle predicts that the number-count
dipole of distant sources is purely kinematic: it arises
from the observer's motion relative to the matter rest
frame, with amplitude $D = [2 + x(1 + \alpha)]\,v/c$,
where $x$ is the number-count slope and $\alpha$ the
spectral index~\cite{EllisBalwin1984}.  For the CMB dipole
velocity $v = 369.82 \pm 0.11$~km/s, the predicted
amplitude is $D_{\rm CMB} \simeq 0.0072$ for mid-infrared
quasars and $D_{\rm CMB} \simeq 0.0043$ for radio
continuum sources.

Three generations of measurements have established that
the observed dipole exceeds this prediction by a factor
of two to four.


\subsection{The CatWISE quasar dipole}
\label{sec:catwise}

Secrest, von Hausegger, Rameez, Mohayaee \&
Sarkar~\cite{Secrest2021} measured the number-count dipole
of $1.36$ million CatWISE2020 quasars at a mean redshift
$\langle z \rangle = 1.2$.  The dipole amplitude is twice
the kinematic expectation, directed $\sim\!28^\circ$ from
the CMB dipole, at a significance of
$4.9\sigma$~\cite{Secrest2021}.  A joint analysis combining
NVSS radio galaxies and CatWISE quasars raised the
rejection to $> 5\sigma$~\cite{Secrest2022}.

The result has been scrutinised for systematics.
Guandalin, Piat, Clarkson \&
Maartens~\cite{Guandalin2023} showed that mask-induced
mode coupling from the $52.6\%$ sky cut can transfer power
from the octupole into the dipole at a level comparable to
the reported excess.  Bashir, Tiwari \&
Zhao~\cite{Bashir2025} performed a comprehensive
reassessment incorporating shot noise, clustering, mask
mode coupling, and cosmic variance via FLASK simulations,
finding a revised significance of
$3.3$--$3.6\sigma$---still anomalous, but more conservative
than the original $4.9\sigma$.  Von Hausegger
et~al.~\cite{vonHausegger2026} reanalysed the CatWISE
multipole structure with a negative binomial likelihood and
found no evidence for octupole contamination, reaffirming the
anomalous dipole.


\subsection{The radio continuum dipole at $5.4\sigma$}
\label{sec:radio-dipole}

B{\"o}hme, Schwarz, Tiwari
et~al.~\cite{Boehme2025} combined three independent radio
surveys---NVSS ($1.4$~GHz), RACS-low ($887.5$~MHz), and
LoTSS-DR2 ($144$~MHz)---and reported a source-count dipole
exceeding the kinematic expectation by a factor of
$3.67 \pm 0.49$, a $5.4\sigma$ discrepancy.  The dipole
direction is within $\sim\!1\sigma$ of the CMB dipole at
$(l, b) \approx (265^\circ, 48^\circ)$.

The methodology introduces a Bayesian estimator based on
the negative binomial distribution, which accounts for
overdispersion caused by multi-component radio sources
(lobes, jets).  Standard Poisson models produce reduced
$\chi^2$ values as high as $23.8$ for TGSS; the negative
binomial model restores $\chi^2 \approx 1$ across all
surveys.  The overdispersion parameter corresponds to
$\sim\!1.1$--$1.3$ mean components per source.

This result is a high-significance radio amplitude excess
reported independently of the mid-infrared CatWISE sample.
The multi-frequency consistency (from $144$~MHz to
$1.4$~GHz) constrains some frequency-dependent systematics, but
does not identify the source of the excess.


\subsection{Cross-wavelength consistency}
\label{sec:cross-wavelength}

Wagenveld, von Hausegger, Kl{\"o}ckner \&
Schwarz~\cite{Wagenveld2025} performed the first joint
Bayesian decomposition of the observed dipole into kinematic
and non-kinematic components, exploiting the different
kinematic scaling of NVSS, RACS-low, and CatWISE.  The key
result is a significant residual non-kinematic dipole with
amplitude $D_{\rm resid} = (0.81 \pm 0.14) \times 10^{-2}$,
offset from the CMB dipole direction by
$39^\circ \pm 8^\circ$.  A common direction between the
kinematic and residual components is disfavoured by model
selection.

Oayda, Land-Strykowski, Lewis \&
Murphy~\cite{OaydaLewis2025} quantified the inter-dataset
tension using suspiciousness statistics: Planck is in
$> 5\sigma$ tension with CatWISE, and NVSS and CatWISE are
concordant at only $0.4\sigma$ despite completely independent
systematics.  The concordance between radio and infrared
strongly suggests a common astrophysical or cosmological
signal.

The Secrest et~al.\ {\it Reviews of Modern Physics}
Colloquium~\cite{SecrestRMP2025} synthesises these results,
reporting a combined significance of $\sim\!6.4\sigma$ via
sample-size-weighted $Z$-scores.  The Colloquium concludes
that the cosmic dipole anomaly ``poses a serious challenge to
FLRW cosmology.''


% ============================================================
\section{Statistical methodology}
\label{sec:dipole-stats}
% ============================================================

The dipole measurements reviewed above employ three
distinct estimators, each with different sensitivity to
systematics.

The \emph{Ellis--Baldwin estimator} computes the dipole
from the variation of source counts across the sky as a
function of flux threshold, using the number-count slope $x$
and spectral index $\alpha$ to predict the kinematic
amplitude.  The estimator is sensitive to the assumed
power-law form of the source counts and to the accuracy of
the spectral index measurement.

The \emph{negative binomial Bayesian estimator} of
B{\"o}hme et~al.~\cite{Boehme2025} replaces the Poisson
likelihood with a negative binomial, accounting for
overdispersion from multi-component sources.  The effect on
the amplitude uncertainty is substantial: Poisson and
negative binomial errors can differ by up to $60\%$.

The \emph{simulation-based inference} approach of Oayda \&
Lewis~\cite{OaydaLewis2026} uses forward-modelled mock
catalogues with survey-specific selection functions to
marginalise over photometric systematics, including the
ecliptic-latitude-dependent Eddington bias induced by the
WISE scanning law.  The CatWISE dipole excess survives this
treatment, though with a $\sim\!3\sigma$ direction
misalignment relative to the CMB dipole that requires
further investigation.

The spread in reported significances---from $3.3\sigma$
(Bashir et~al.) to $6.4\sigma$ (Secrest
et~al.\ combined)---reflects the dependence on the
estimator and the treatment of systematics, not a
fundamental disagreement about the existence of the excess.
All analyses agree that the observed dipole exceeds the
kinematic prediction; the debate concerns the precise
amplitude and significance after accounting for
mode coupling, overdispersion, and photometric errors.


% ============================================================
\section{Peculiar velocity surveys and bulk flow}
\label{sec:dipole-pv}
% ============================================================

The number-count dipole constrains the matter rest frame
through source statistics.  Peculiar velocity surveys
constrain the same rest frame through direct distance
measurements, providing an independent cross-check.


\subsection{CosmicFlows-4 and CF4++}
\label{sec:cf4}

Watkins et~al.~\cite{Watkins2023} measured a bulk flow of
$419 \pm 36$~km/s at $200\,h^{-1}$~Mpc from the
CosmicFlows-4 catalogue ($55{,}877$ galaxies), with a
$\Lambda$CDM probability below $0.003\%$.  Courtois
et~al.~\cite{Courtois2025} extended the sample to CF4++
($65{,}518$ galaxies, incorporating WALLABY, FAST/FASHI,
and DESI-PV data) and found $315 \pm 40$~km/s at
$150\,h^{-1}$~Mpc ($1.7\sigma$ tension) and
$351 \pm 109$~km/s at $200\,h^{-1}$~Mpc ($2.5\sigma$
tension).  The dynamical homogeneity scale---the depth at
which the bulk flow converges to zero---is not reached
within the $200$--$300\,h^{-1}$~Mpc interval.


\subsection{DESI DR1 peculiar velocities}
\label{sec:desi-pv}

The DESI DR1 peculiar velocity programme~\cite{DESIDR1PV2025}
delivered $76{,}616$ measurements---the largest peculiar
velocity sample ever assembled---using both Fundamental Plane
($\sim\!66{,}000$ early-type galaxies) and Tully--Fisher
($\sim\!10{,}000$ spirals) distance indicators.  Three
independent analysis methods yield a consensus growth rate
$f\sigma_8(z_{\rm eff} = 0.07) = 0.450 \pm 0.055$, in
remarkable agreement with the Planck+$\Lambda$CDM prediction
of $0.449 \pm 0.008$.  The growth index
$\gamma = 0.58 \pm 0.11$ is consistent with the GR prediction
$\gamma_{\rm GR} = 0.55$.

This result creates a tension within the anomaly picture that
we must address honestly.  The DESI growth-rate measurement
finds no departure from $\Lambda$CDM, yet the CF4++ bulk flow
remains anomalously large at $> 150\,h^{-1}$~Mpc.  The
decoupling between the growth rate (which measures the rate of
structure formation through velocity-density correlations) and
the bulk flow amplitude (which measures the large-scale
coherent motion) suggests that the anomaly is not driven by
anomalous gravitational collapse but by a large-scale
kinematic mode that does not affect the local growth rate.
A cosmological tilt in the Bianchi framework is precisely such
a mode: it contributes a coherent bulk velocity without
modifying the local gravitational dynamics.


% ============================================================
\section{Theoretical status of candidate mechanisms}
\label{sec:dipole-theory}
% ============================================================

Published measurements report that the matter dipole exceeds
the FLRW kinematic prediction.  We now
assess which theoretical mechanisms can produce this excess
within the constraints imposed by the CMB, nucleosynthesis,
and gravitational-wave observations.  A detailed discussion
of each mechanism is deferred to
Chapter~\ref{ch:discussion}
(Section~\ref{sec:disc-mechanisms}); here we state the
conclusions.


\paragraph{Mechanisms ruled out.}

Three developments since 2025 have eliminated previously
viable candidates.
Clarkson \& Maartens~\cite{ClarksonMaartens2026}
demonstrated that the enhanced peculiar-velocity growth
claimed by Tsagas and collaborators ($v \propto a^{3/2}$)
arises from an inconsistent treatment of constraint
equations in the 1+3 covariant formalism; when constraints
are imposed consistently, the standard growth law
$v \propto a^{1/2}$ is recovered.
Mart{\'i}n, Skordis et~al.~\cite{MartinSkordis2025}
systematically evaluated four tilt sources in Bianchi~V and
found that spatial curvature, heat flow, and electromagnetic
fields each fall short of the observed dipole amplitude by
five to thirteen orders of magnitude.


\paragraph{Surviving mechanisms.}

Two mechanisms remain viable within the tilted-Bianchi
framework.  The khronon field---a preferred-frame scalar
field that is the hypersurface-orthogonal limit of
Einstein-aether theory---achieves
$\beta \sim 2$--$4 \times 10^{-3}$ within the
post-GW170817 parameter
space~\cite{MartinSkordis2025}.
The Krishnan--Mondol--Sheikh-Jabbari tilt
instability~\cite{Krishnan2023}, which operates at the exact
(non-perturbative) Bianchi~V level, shows that the cosmic
no-hair theorem forces metric shear $\sigma \to 0$ but does
not force the fluid tilt $\beta \to 0$---tilt can persist
or grow during $\Lambda$-dominated expansion.

Outside the Bianchi framework, superhorizon axion-induced
isocurvature perturbations offer a competing
explanation that generates an intrinsic negative CMB dipole
partially cancelling the kinematic component.  Because the
corresponding quadrupole prediction is not yet attached here to
a verified citation or generated source artifact, it is treated
as external motivation rather than a current manuscript result.


\paragraph{Separation principle.}

The departure framework is mechanism-agnostic by construction:
the master identity (Eq.~\ref{eq:master-identity}) is an
algebraic relation among geometry-frame kinematic quantities
and does not assume a specific tilt source.  Amplitude fitting
and mechanism attribution are separable questions: the current
scalar pipeline fits an admitted beta-like amplitude but does
not establish that its source is cosmological; the theoretical
analysis constrains candidate mechanisms
(Section~\ref{sec:disc-separation}).


% ============================================================
\section{Observer-frame discriminator: separating kinematic dipole from Bianchi tilt}
\label{sec:dipole-discriminator}
% ============================================================

The operational question raised by the dipole anomaly is not
merely whether a dipole exists, but whether the observed sky is
better described as a local observer boost or as a genuine
cosmological tilt.  The FB-8 observer-frame layer makes that
question explicit by leaving the FB-7 likelihood in the
cosmological frame and adding a separate observer-frame wrapper
with a type-distinct `ObserverBoost` carrier.  The resulting
workflow enforces the physically relevant order
``cosmo-tilt $\rightarrow$ observer-boost`` rather than
collapsing both effects into a single effective velocity.


\subsection{Linear aberration kernel}

For an observer moving at speed $\beta_{\rm obs}$ relative to
the cosmological frame, the temperature and polarisation
multipoles are mixed by a linear aberration kernel
$K_{\ell\ell'}(\beta_{\rm obs})$ of the
Challinor--van Leeuwen form~\cite{ChallinorVanLeeuwen2002}.  In
the aligned case used by the present implementation, the kernel
is tridiagonal at ${\cal O}(\beta_{\rm obs})$: power leaks only
between neighbouring multipoles, while the
$\beta_{\rm obs} = 0$ limit returns the identity exactly.  The
physically relevant scale is the Solar-system dipole speed
$\beta_{\rm obs} \simeq 1.23\times 10^{-3}$ identified by the
Planck observer-velocity analysis~\cite{Planck2013XXVII}, so the
observer-frame correction is small but not negligible at the
low-$\ell$ level where the dipole anomaly is analysed.


\subsection{Observer-frame $C_\ell$ and $a_{\ell m}$ adapters}

The kernel acts on both the harmonic coefficients and the
diagonal spectra.  The $a_{\ell m}$ adapter applies the aligned
mixing directly in multipole space, whereas the $C_\ell$
adapter propagates the same linear leakage into the diagonal
power spectra.  These two views are intentionally kept as
separate audited surfaces because they answer different
questions: the harmonic map shows the neighbouring-$\ell$
transfer explicitly, while the spectrum adapter isolates the
observable change in $C_\ell^{TT}$, $C_\ell^{EE}$,
$C_\ell^{TE}$, and $C_\ell^{BB}$.  A genuine cosmological tilt
changes the rest-frame sky before any observer correction is
applied; a local boost instead remaps a pre-existing
cosmological-frame sky.  The distinction becomes visible once
both adapters are pinned to the same linear kernel.


\subsection{Composition order and non-commutation}

Observer motion and cosmological tilt do not commute.  If one
first constructs a cosmological-frame tilted sky and then boosts
to the local frame, the resulting multipoles differ from those
obtained by first applying an observer boost to an untilted sky
and only then trying to absorb the result into a global tilt.
This is not a bookkeeping nuisance but the key physical point:
observer aberration is a local frame transformation, whereas a
Bianchi tilt changes the cosmological-frame source itself.  The
implementation therefore keeps \texttt{compose\_tilts(...)} as a
diagnostic witness only, never as a production shortcut, and
pins the ordering
``cosmo-tilt $\rightarrow$ observer-boost'' in the conventions
document and the audit ledger.


\subsection{Likelihood-ratio discriminator and coverage}

The operational discriminator is the likelihood-ratio statistic
\begin{equation}
\Lambda(d; H_{\rm obs}, H_{\rm cosmo})
= 2\,[\ln L_{\rm cosmo}(d) - \ln L_{\rm obs}(d)]\,.
\end{equation}
Negative $\Lambda$ favours the local-boost explanation
$H_{\rm obs}$; positive $\Lambda$ favours the cosmological-tilt
explanation $H_{\rm cosmo}$.  The observer-frame likelihood is
constructed by wrapping the FB-7 cosmological-frame likelihood
with an explicit observer-boost prior rather than reopening the
cosmological-frame class.  This preserves the type-distinct
surface and makes the model comparison auditable.  Following the
recoverability argument of Kosowsky \&
Kahniashvili~\cite{KosowskyKahniashvili2011}, the phase-close
verification uses $N=500$ seeded synthetic draws under each
hypothesis and checks that the empirical PIT p-values remain
consistent with a uniform distribution at the $99\%$ level.
the conditioned legacy figure gallery in Appendix~\ref{app:conditioned-legacy-figure-gallery} shows the resulting
coverage histograms for both truth hypotheses.


% ============================================================
\section{Near-term experimental outlook}
\label{sec:dipole-outlook}
% ============================================================

Three experiments will transform the anomaly picture within
the next two to five years.

\emph{Euclid DR1} (scheduled October~2026) will provide
$\sim\!2100~\mathrm{deg}^2$ of wide-survey photometric data,
enabling the first independent large-area Ellis--Baldwin
dipole test at optical/near-infrared wavelengths.  The Euclid
galaxy sample ($\langle z \rangle \sim 0.9$) is
systematically independent of both CatWISE (mid-infrared) and
NVSS/RACS (radio), making it the most powerful single-dataset
test of the anomaly in the near term.

\emph{Simons Observatory} first-light data
($\sim\!2027$) will deliver CMB polarisation maps with noise
levels $\sim\!3\times$ below Planck, enabling BiPoSH analyses
and hemispherical asymmetry tests in polarisation at
unprecedented sensitivity.

\emph{SKA Phase~1} ($\sim\!2028$--$2030$) will survey the
radio continuum dipole with a factor-of-ten improvement in
source density over existing surveys, pushing the radio-only
significance beyond $10\sigma$ if the anomaly is real.  The
EMU precursor survey ($\sim\!20$ million expected
extragalactic sources) will provide an interim test.

The eROSITA X-ray catalogue ($\sim\!930{,}000$ sources from
eRASS1) could fill a gap in the multi-wavelength picture, but
no dedicated X-ray cosmic dipole measurement has been
published as of March~2026.


% ============================================================
\section{Anomaly parameters for the departure framework}
\label{sec:dipole-summary}
\label{sec:reframing}
% ============================================================

We summarise the observational inputs that enter the departure
framework of Chapter~\ref{ch:framework}.

The tilt rapidity $\beta_{\rm CF4} = 1.334 \times 10^{-3}$
is derived from the Watkins et~al.~\cite{Watkins2023}
CosmicFlows-4 bulk flow of $400 \pm 30$~km/s at
$200\,h^{-1}$~Mpc, converted to a rapidity via
$\beta = v/c$.  This is the \emph{input} value from the
CF4 catalogue; the VER06 posterior median after
multi-channel inference is
$\beta = 1.360 \times 10^{-3}$
(Chapter~\ref{ch:results}), differing by $2\%$ due to
the pull of the CMB and number-count channels.
The dipole direction is consistent across
radio, infrared, and peculiar velocity surveys within the
directional coherence tolerance, at
$(l, b) \approx (264^\circ, 48^\circ)$.

The number-count dipole significance ranges from
$\sim\!3.3\sigma$ (conservative, incorporating mode-coupling
and photometric uncertainty;
Bashir et~al.~\cite{Bashir2025}) to $\sim\!6.4\sigma$
(combined $Z$-scores across all surveys;
Secrest et~al.~\cite{SecrestRMP2025}).  The departure
framework records the legacy HTT posterior mean
$\bar{Q}\simeq 0.092$ (95\% HPD $[0.035,0.19]$) for the
sign-clean flat-comparator FLRW$_{\mathrm{tilt}}$ run, and
separately records the S3 dust filling fraction
$\FF=0.063^{+0.029}_{-0.023}$ (Chapter~\ref{ch:results}).
These transfer-conditional summaries integrate over the full
$\beta$ posterior rather than relying on a single-point
significance estimate, but they are not MIO diagnostic
certificates and do not carry geometry or family-ID semantics.

The non-kinematic residual dipole identified by Wagenveld
et~al.~\cite{Wagenveld2025}---offset by $39^\circ \pm 8^\circ$
from the CMB dipole---is not incorporated into the VER06
pipeline, which assumes a common dipole direction for all
channels.  Incorporating the directional offset is a target
for the direction-dependent programme
(Section~\ref{sec:phase30}).

\begin{remark}[EGS theorem and isotropy]
\label{rem:EGS}
The Ehlers--Geren--Sachs theorem establishes that if all
fundamental observers in an expanding universe measure an
exactly isotropic CMB, the spacetime is exactly FLRW.
Almost-EGS extensions allow for small anisotropies
consistent with observations.  The dipole anomaly does
not violate the EGS theorem directly---the CMB is
isotropic to $10^{-5}$ in the CMB rest frame---but it
implies that the matter rest frame is displaced from the
CMB rest frame by more than the $\Lambda$CDM kinematic
prediction.  The departure framework quantifies this
displacement through the tilt parameter $\Otilt$.
\end{remark}
